\section{The constant produced by the proof of the area law}
\label{sec:area_law_explicit_constant}

The proof of Theorem~1.1 of \cite{OpenAI2026AreaLaw} produces the constant of
the area law from the constants of Propositions~9.6, 10.2 and~11.2. This
section records that constant.

% Status: these are auxiliary results. The source asserts only that a constant
% exists; the statement of Theorem 1.1 is thm:openai_uniform_area_law.

% Source: 10-geometry.tex, proof of Theorem 1.1, lines 803--857.
% Revision: adc7f1241b42e322a6451854ab7e4b4c146bf78a.
\begin{definition}[The constant in the two-dimensional area-law proof]
    \label{def:area_law_proof_constant}
    \lean{TNLean.PEPS.AreaLaw.uniformAreaLawConstantOf,
        TNLean.PEPS.AreaLaw.uniformAreaLawScale}
    \leanok
    For an integer $q\ge1$ and real numbers $C_2,C_3$ put $c_D=\max\{C_3,0\}$,
    $c_M=\max\{C_2,0\}$ and
    \begin{align}
        C(q,C_2,C_3) & =c_D\log q+\frac{c_Mc_D}{2}.
    \end{align}
    For integers $N_1,N_2$ and a real number $C_2$, let $N_3$ be the collar
    radius threshold at the coefficient $C_2$ and the factor $1$, and put
    $n_*(N_1,N_2,C_2)=\max\{N_1,N_2,N_3\}$.
\end{definition}

\begin{theorem}[The area law at given constants of the three propositions]
    \label{thm:area_law_conditional_assembly_constants}
    \lean{TNLean.PEPS.AreaLaw.regionalEntropy_le_of_collar_amplification_twoFamilies_constants,
        TNLean.PEPS.AreaLaw.uniformAreaLawConstantOf_nonneg,
        TNLean.PEPS.AreaLaw.two_le_uniformAreaLawScale}
    \leanok
    \uses{def:area_law_proof_constant,def:area_law_collar_bound,
        def:area_law_amplification_bound,def:area_law_two_family_decomposition,
        def:area_law_ground_entropy}
    Fix $q\ge1$, $R\ge0$, $J$, $\Delta>0$ and template constants
    $D_0,C_{\rm tpl}$. Suppose that the assertions of the sublinear collar
    bound hold with the constant $C_1$ for all $n\ge N_1$; that the assertions
    of the amplification bound for the input constant $\max\{C_1,1\}$ and the
    order $M=100$ hold with the constant $C_2$ for all $n\ge\max\{2,N_2\}$;
    and that $C_3$ is a constant of the two-family decomposition at the lower
    scale $n_*(N_1,N_2,C_2)$. Then $C(q,C_2,C_3)\ge0$, $n_*(N_1,N_2,C_2)\ge2$,
    and for every finite induced domain $\Lambda$, every admissible
    Hamiltonian with a gapped ground vector $\Omega$ and every
    $A\subseteq\Lambda$,
    \begin{align}
        S_\Omega(A) & \le C(q,C_2,C_3)\,|\partial_\Lambda A|.
    \end{align}
\end{theorem}
\begin{proof}
    \leanok
    \uses{thm:area_law_zero_boundary,
        thm:area_law_boundary_assembly,thm:area_law_uniform_collar_threshold}
    The argument is the proof of Theorem~1.1 of \cite{OpenAI2026AreaLaw} with
    the constants fixed. If $|\partial_\Lambda A|=0$ the entropy vanishes by
    Theorem~\ref{thm:area_law_zero_boundary}. The collar radius threshold $N_3\ge2$ gives
    $L_n+2r<n^\beta$ for every $n\ge N_3$ and every $r\le C_2n^{(1-\varepsilon)/\alpha}$.
    Each piece has scale $n_i\ge n_*\ge\max\{2,N_1,N_2,N_3\}$, so the collar
    and amplification bounds apply to it and give
    $I_\Omega(X_i:J_i)\le c_Mn_i^{-100}$. The boundary estimate from partition
    bounds, with $|D|\le c_Db$ and $\sum_ic_Mn_i^{-100}\le c_Mc_Db$, gives the
    constant $c_D\log q+c_Mc_D/2$.
\end{proof}

% Source: 10-geometry.tex, proof of Theorem 1.1, lines 803--857, at the
% template constants D_0 = 2R + 11 and C_tpl = 24.
% Revision: adc7f1241b42e322a6451854ab7e4b4c146bf78a.
\begin{definition}[The constant of the area law from the named constants]
    \label{def:area_law_named_constant}
    \lean{TNLean.PEPS.AreaLaw.uniformAreaLawConstant}
    \leanok
    \uses{def:area_law_proof_constant,def:area_law_two_family_explicit_constant}
    Fix $q\ge1$, $R\ge0$, $J$ and $\Delta$. Let $C_{\rm box}$ be a safe-box
    constant at safety parameter $2R+11$, let
    $C_{\rm tr},c_t,C_{\rm ent},e_{\rm ent},C_{\rm en},e_{\rm en}$ be
    truncation and transport constants of the scanner, let
    $C_{\rm loc},c_{\rm loc}$ be filter-locality constants of the positive
    constraints, and put $c_1=1/112$. Let $C_1$ and
    $N_1$ be the constant and the threshold of the sublinear collar bound at
    these data, put $C_0=\max\{C_1,1\}$, let $C_2$ and $N_2'$ be the constant
    and the threshold of the amplification bound at $C_0$ and the order
    $M=100$, and let $N_2=\max\{N_2',N_3\}$ with $N_3$ the collar radius
    threshold at $C_2$ and the factor $1$. With
    $n_*=n_*(N_1,N_2,C_2)$ and the two-family constant
    $C_3=C_{\rm expl}(2R+11,24,n_*,c_1)$, the constant of the area law is
    $C(q,C_2,C_3)$.
\end{definition}

\begin{theorem}[The area law with the named constant, from its inputs]
    \label{thm:area_law_named_constant}
    \lean{TNLean.PEPS.AreaLaw.exists_regionalEntropy_le_uniformAreaLawConstant}
    \leanok
    \uses{def:area_law_named_constant,def:area_law_box_entropy,
        def:area_law_ground_entropy}
    Fix $q\ge1$, $R\ge0$, $J\ge0$ and $\Delta>0$. Assume that
    $C_{\rm box}\ge0$ satisfies $F_{\rm box}(r)\le C_{\rm box}r^{1+e_*}$ at
    safety parameter $2R+11$ for every $r\ge1$; that $C_{\rm loc}\ge0$ and
    $c_{\rm loc}>0$ are filter-locality constants of the positive constraints
    used in the proof of Proposition~10.2. Then there are truncation and transport
    constants of the scanner, with $C_{\rm tr},c_t>0$, such that the constant
    $C$ of Definition~\ref{def:area_law_named_constant} at these data is
    nonnegative, and for every finite induced domain $\Lambda$, every
    admissible Hamiltonian with a gapped ground vector $\Omega$ and every
    $A\subseteq\Lambda$,
    \begin{align}
        S_\Omega(A) & \le C\,|\partial_\Lambda A|.
    \end{align}
\end{theorem}
\begin{proof}
    \leanok
    \uses{thm:area_law_conditional_assembly_constants,
        thm:area_law_sublinear_collar,thm:area_law_amplification,
        thm:area_law_two_family_explicit_constant}
    The safe-box bound gives truncation and transport constants for which the
    sublinear collar bound holds with the constant $C_1$ from the threshold
    $N_1$ on, at $D_0=2R+11>2R+10$ and $C_{\rm tpl}=24$. The filter-locality
    constants give the amplification bound at the input constant $C_0$ and
    the order $100$ with the constant $C_2$ from the threshold
    $\max\{N_2',N_3\}$ on. Since $1/112$ is a birth-separation constant,
    $C_{\rm expl}(2R+11,24,n_*,1/112)$ is a constant of the two-family
    decomposition at $n_*$. Theorem~\ref{thm:area_law_conditional_assembly_constants}
    then gives the bound.
\end{proof}
